\section*{Abstract}
Multimodal models trained by empirical risk minimization can exploit spurious visual correlations in their training data instead of the physiological signal they are intended to use. This thesis examines this risk for stress detection in two lines of work. In the first, a pre-registered test on four participants from the UBFC-Phys dataset shows that facial video itself contains recoverable physiological information (POS estimator, Mann--Whitney $p = 0.0314$; exploratory CHROM $p = 0.0001$ and PBV $p = 0.0103$), so a visual branch can have access to genuine as well as spurious cues. In the second, the thesis introduces EQUITAS-RCMF, a multimodal architecture that projects MobileNetV3-Small visual features onto orthogonal causal and confounder subspaces with a thin QR projection, uses a scar mask as privileged information during training only, and gates visual evidence according to how strongly it concentrates on the scar. EQUITAS-RCMF is evaluated on a controlled benchmark in which chest physiology from 15 WESAD participants is paired at random with 418 FFHQ face images carrying a synthetic brow-line scar, so that the scar is the only label-related visual cue; in training, the scar accompanies stress in 85\% of samples. On data from three held-out participants, the model's accuracy in its deployed, mask-free mode is 61.7\% whether the scar is positively associated with stress, independent of it, or negatively associated with it; removing the scar changes its predictions by about 0.0006 on average, and its equalized-odds gap stays at or below 0.065. The compiled model runs at 0.90~ms per frame on a desktop CPU. The thesis also discusses the limits of these results, including the small cohorts, a physiology-only comparison on the same data that is still pending, and a difference in scale between the focus signal used in training and at inference, and it examines the design's implications for data protection and algorithmic accountability.

\vspace{0.75cm}

\noindent\textbf{Keywords:} Counterfactual Fairness; Shortcut Learning; Edge AI; Stiefel Manifold; Learning Using Privileged Information; MobileNetV3; Remote Photoplethysmography; Stress Detection
